\documentclass[11pt,a4paper]{article}
\usepackage[T1]{fontenc}
\usepackage[utf8]{inputenc}
\usepackage{lmodern}
\usepackage{amsmath,amssymb,amsthm}
\usepackage{booktabs}
\usepackage{xcolor}
\usepackage[margin=27mm]{geometry}
\usepackage[colorlinks=true,urlcolor=blue,citecolor=blue,linkcolor=blue]{hyperref}
\newtheorem{theorem}{Theorem}
\newtheorem{proposition}[theorem]{Proposition}
\newtheorem{lemma}[theorem]{Lemma}
\newtheorem{corollary}[theorem]{Corollary}
\newtheorem{conjecture}[theorem]{Conjecture}
\theoremstyle{definition}
\newtheorem{definition}[theorem]{Definition}
\newtheorem{remark}[theorem]{Remark}
\theoremstyle{plain}
\DeclareMathOperator{\Tr}{Tr}
\DeclareMathOperator{\Spec}{Spec}
\newcommand{\ket}[1]{\lvert#1\rangle}
\newcommand{\bra}[1]{\langle#1\rvert}
\newcommand{\braket}[2]{\langle#1\vert#2\rangle}
\newcommand{\norm}[1]{\lVert#1\rVert}
\newcommand{\one}{\mathbf{1}}
\newcommand{\C}{\mathbb{C}}
\newcommand{\Nop}{\hat N}
\setlength{\emergencystretch}{2em}
\setlength{\parindent}{0pt}
\setlength{\parskip}{0.3em plus 0.1em minus 0.1em}
\title{Spectra of particle-reduced fermionic density operators:\\
the Amosov--Filippov conjecture holds up to a factor $p!$}
\author{Fable \& Astra\\[0.35em]
\normalsize Research led by Jean-Philippe Bruneton\\[-0.1em]
\small Universit\'e Paris Cit\'e, France}
\date{19 November 2026}

\begin{document}
\maketitle
\begin{abstract}
The $p$-particle reduced density matrix of a many-fermion state keeps exactly
the information needed to evaluate $p$-body observables, and for a pure state
with a fixed number $N$ of particles a classical duality, usually attributed
to Carlson and Keller, says that keeping $p$ particles and keeping $N-p$
particles produce the same nonzero spectrum. Amosov and Filippov asked what replaces this duality on
the whole fermionic Fock space, where the particle number need not be fixed.
They introduced the map $\Phi(\varrho)=\sum_j a_j\varrho a_j^\dagger$, which
removes one particle, and conjectured that, for every pure state, the operator
$\Phi^p(\varrho)$ obtained by removing $p$ particles has the same spectrum as
the $p$-particle reduced density operator $\varrho_p$. We settle the
conjecture. For every pure state $\varrho$ on the Fock space of $n$ modes and
every $p$, the full spectra satisfy
$\Spec\Phi^p(\varrho)=p!\,\Spec\varrho_p$, zero eigenvalues included. The
conjecture is therefore true as stated for $p=1$ and false for every $p\ge2$
unless both operators vanish; the minimal counterexample is the two-mode state
$\ket{11}$. The proof identifies $\varrho_p$ and $\Phi^p(\varrho)/p!$ as the
two complementary outputs of the completely positive map with Kraus operators
$a_{s_p}\cdots a_{s_1}$. The same factor corrects the fixed-particle-number
identity of the source, which becomes $\Phi^{N-p}(\varrho)=(N-p)!\,\varrho_p$.
\end{abstract}

\section{Introduction}
A state of $n$ fermionic modes is a density operator on the Fock space $H_n$,
of dimension $2^n$. Many physical quantities depend on such a state only
through a few particles at a time: the kinetic energy is a one-body
observable, a pair interaction is a two-body observable. The object that
carries exactly this information is the $p$-particle reduced density matrix,
the central object of reduced-density-matrix methods in quantum chemistry.
For a pure state of exactly $N$ fermions, the $p$-particle and
$(N-p)$-particle reduced density matrices have the same nonzero eigenvalues, a
classical result that Amosov and Filippov attribute to Carlson and
Keller~\cite{AmosovFilippov2017} and that Gigena, Di Tullio and Rossignoli
derive from a singular value decomposition~\cite[Sec.~II\,B]{Gigena2021}. This is the fermionic analogue
of the familiar fact that the two halves of a bipartite pure state have the
same spectrum: removing $N-p$ particles and keeping $p$ of them are
spectrally equivalent operations.

States of fermions need not have a fixed particle number. Superconducting and
other quasi-free states mix particle-number sectors, and the parity
superselection rule only forbids superpositions of even and odd numbers.
Amosov and Filippov \cite{AmosovFilippov2017} looked for the analogue of the
Carlson--Keller duality on the whole Fock space. They introduced the linear
map $\Phi(\varrho)=\sum_{j=1}^n a_j\varrho a_j^\dagger$, which removes one
particle from any operator, proved that $\Phi^{N-p}(\varrho)$ equals the
$p$-particle reduced density operator $\varrho_p$ when the particle number is
fixed, and conjectured that for every pure state on $H_n$, with no assumption
on particle number or parity, $\Phi^p(\varrho)$ and $\varrho_p$ have the same
spectrum. They checked the case $p=1$, $n=2$ by hand and reported that other
examples, general pure states of three and four modes, are also
equispectral.

The question answered here is whether this conjecture holds, and what the
correct statement is if it does not. The answer is a single identity. For
every pure state and every $p$,
\begin{equation}\label{eq:main}
\Spec\Phi^p(\varrho)=p!\;\Spec\varrho_p ,
\end{equation}
for the full spectra of the two operators on $H_n$, counted with
multiplicity. The conjecture is thus true for $p=1$ and false for every
$p\ge2$ except in the trivial case where both operators vanish. The factor
$p!$ counts the orderings in which $\Phi^p$ removes the same set of $p$
particles, and it is already visible in the traces: $\Tr\varrho_p$ is the
expected number of $p$-element subsets of occupied modes, whereas
$\Tr\Phi^p(\varrho)$ is the expected number of ordered $p$-tuples. The
simplest example shows everything. For two modes and the state $\ket{11}$ in
which both are occupied, $\varrho_2=\ket{11}\bra{11}$ has spectrum
$\{1,0,0,0\}$, while removing both particles one after the other in either
order gives $\Phi^2(\varrho)=2\ket{00}\bra{00}$, with spectrum $\{2,0,0,0\}$.
This state has a fixed particle number $N=2$, so the corollary that the
source derives for fixed particle number also requires the factor.

Behind \eqref{eq:main} lies a structural fact. The $p$-particle reduced
density operator and $\Phi^p(\varrho)/p!$ are the two complementary outputs,
in the sense of a Stinespring dilation, of one completely positive map whose
Kraus operators are the ordered products of $p$ annihilation operators. Two
complementary outputs of a completely positive map acting on a pure input
always share their nonzero spectrum, exactly as the two marginals of a
bipartite pure state do. The particle-number-free extension sought by Amosov
and Filippov therefore exists, with $\Phi^p(\varrho)/p!$ in the role of the
$(N-p)$-particle reduced density operator when $N$ is no longer defined.

Section~\ref{sec:related} relates the result to the literature.
Section~\ref{sec:setting} recalls the definitions of the source.
Section~\ref{sec:results} proves the main theorem and its consequences for
the conjecture and for the fixed-particle-number proposition of the source.
Section~\ref{sec:verification} describes the numerical checks and
Section~\ref{sec:scope} the scope of the result.

\section{Related work}\label{sec:related}
The spectral duality between the $p$-particle and $(N-p)$-particle reduced
density matrices of a pure $N$-fermion state is classical; Amosov and
Filippov attribute it to Carlson and Keller and use it to derive the
corollary of their Proposition~4~\cite{AmosovFilippov2017}. It is the
fixed-$N$ special case of Theorem~\ref{thm:main} combined with
Corollary~\ref{cor:fixedN} below.

Amosov and Filippov \cite{AmosovFilippov2017} define the reduced density
operator $\varrho_p$ on the full Fock space, introduce the particle-removing
map $\Phi$, and study spectral properties of reductions under the parity
superselection rule. Their Proposition~4 states
$\Phi^{N-p}(\varrho)=\varrho_p$ for states with fixed particle number $N$,
and the corollary that follows it deduces equality of the spectra of
$\varrho_p$ and $\Phi^p(\varrho)$ for pure fixed-$N$ states. The conjecture
treated here closes their Section~5. We prove that Proposition~4 holds
exactly for $N-p\le1$ and requires the factor $(N-p)!$ otherwise. Their proof
correctly shows that $\Phi$ maps an $m$-particle operator to its reduction by
one particle; iterating this relation counts each ordered removal separately,
which produces the factorial.

Gigena, Di Tullio and Rossignoli \cite{Gigena2021} study many-body
entanglement in fermion systems through $M$-body density operators. For an
$N$-fermion state they write
$\hat\rho^{(M)}=\frac{1}{(N-M)!}\sum_{j_1,\dots,j_{N-M}}c_{j_1}\cdots
c_{j_{N-M}}\,\rho\,c^\dagger_{j_{N-M}}\cdots c^\dagger_{j_1}$
\cite[Eq.~(79)]{Gigena2021}, normalized so that $\Tr\rho^{(M)}=\binom NM$
\cite[Eq.~(16)]{Gigena2021}, and they show that, for a pure state,
$\rho^{(M)}$ and $\rho^{(N-M)}$ have the same spectrum apart from the number
of zero eigenvalues \cite[Sec.~II\,B]{Gigena2021}. Their prefactor
$1/(N-M)!$ is exactly the correction of Proposition~4 established in
Corollary~\ref{cor:fixedN}; they refer to Amosov and Filippov for the
extension to states without fixed particle number, which they do not treat,
and they do not comment on the discrepancy with Proposition~4. The present
note supplies that extension.

\section{Setting and the conjecture}\label{sec:setting}
The Fock space $H_n$ of $n$ fermionic modes has the occupation basis
$\ket{j_1\dots j_n}$, $j_s\in\{0,1\}$. Following \cite{AmosovFilippov2017},
the creation operators act by
\[
a_t^\dagger\ket{j_1\dots j_n}=(-1)^{j_1+\cdots+j_{t-1}}
\ket{j_1\dots j_{t-1}\,1\,j_{t+1}\dots j_n}\quad\text{if } j_t=0,
\]
and $a_t^\dagger\ket{j_1\dots j_n}=0$ if $j_t=1$; the annihilation operator
$a_t$ is the adjoint of $a_t^\dagger$. These operators satisfy the canonical
anticommutation relations. The particle-number operator is
$\Nop=\sum_s a_s^\dagger a_s$. For a density operator $\varrho$ on $H_n$ we
write $\langle X\rangle=\Tr(\varrho X)$.

For a subset $S=\{s_1<\dots<s_p\}$ of $\{1,\dots,n\}$ we set
$A_S=a_{s_p}\cdots a_{s_1}$, so that
$A_S^\dagger=a_{s_1}^\dagger\cdots a_{s_p}^\dagger$, and
$\ket{S}=A_S^\dagger\ket{0}$, where $\ket0$ is the vacuum. With the sign
convention above, $\ket S$ is the occupation vector whose occupied modes are
those of $S$, with sign $+1$, because each $a_{s_q}^\dagger$ acts on a vector
whose modes of index below $s_q$ are empty. The vectors $\ket S$ with $|S|=p$
form an orthonormal basis of the $p$-particle sector.

\begin{definition}[\cite{AmosovFilippov2017}]\label{def:rdm}
The $p$-particle reduced density operator of $\varrho$ is the operator on
$H_n$
\[
\varrho_p=\sum_{\substack{s_1<\dots<s_p\\ t_1<\dots<t_p}}
\langle a_{s_1}^\dagger\cdots a_{s_p}^\dagger a_{t_p}\cdots a_{t_1}\rangle\;
a_{t_1}^\dagger\cdots a_{t_p}^\dagger a_{s_p}\cdots a_{s_1}
\prod_{s\notin\{s_1,\dots,s_p,t_1,\dots,t_p\}}a_sa_s^\dagger .
\]
In the notation above,
$\varrho_p=\sum_{|S|=|T|=p}\langle A_S^\dagger A_T\rangle\,E_{T,S}$ with
$E_{T,S}=A_T^\dagger A_S\prod_{s\notin S\cup T}a_sa_s^\dagger$. No
normalization is applied, so $\Tr\varrho_p$ need not equal $1$.
\end{definition}

The particle-removing map is $\Phi(X)=\sum_{j=1}^n a_jXa_j^\dagger$, and
$\Phi^p$ denotes its $p$-fold composition.

\begin{conjecture}[Amosov--Filippov \cite{AmosovFilippov2017}]
For every pure state $\varrho=\ket\psi\bra\psi$ on $H_n$, with no assumption
on particle number or parity, the operators $\varrho_p$ and $\Phi^p(\varrho)$
have the same spectrum.
\end{conjecture}

\section{Results}\label{sec:results}
\subsection{Two lemmas}
The first lemma shows that the product of projectors in
Definition~\ref{def:rdm} turns each term into a matrix unit of the
$p$-particle sector.

\begin{lemma}\label{lem:units}
For $|S|=|T|=p$, $E_{T,S}=\ket T\bra S$. Consequently $\varrho_p$ is
supported on the $p$-particle sector, where
$\bra T\varrho_p\ket S=\Tr(\varrho A_S^\dagger A_T)$, and it vanishes on all
other sectors.
\end{lemma}

\begin{proof}
Let $\ket J$ be an occupation vector with set of occupied modes $J$. The
product $\prod_{s\notin S\cup T}a_sa_s^\dagger$ is the projector onto vectors
whose modes outside $S\cup T$ are empty, so it fixes $\ket J$ if
$J\subseteq S\cup T$ and annihilates it otherwise. Next, $A_S\ket J$ vanishes
unless $S\subseteq J$, in which case it equals $\pm\ket{J\setminus S}$.
Finally, $A_T^\dagger\ket{J\setminus S}$ vanishes as soon as $T$ meets
$J\setminus S$; since $J\setminus S\subseteq T$, only $J=S$ survives. For
$J=S$ we have $A_S\ket S=A_SA_S^\dagger\ket0=\ket0$, because each factor
$a_sa_s^\dagger$ that appears acts on a vector in which mode $s$ is empty, and
$A_T^\dagger\ket0=\ket T$. Hence $E_{T,S}\ket J=\delta_{J,S}\ket T$.
\end{proof}

\begin{lemma}\label{lem:phip}
For every operator $X$ on $H_n$ and $1\le p\le n$,
$\Phi^p(X)=p!\sum_{|T|=p}A_TXA_T^\dagger$.
\end{lemma}

\begin{proof}
By definition
$\Phi^p(X)=\sum_{(j_1,\dots,j_p)}a_{j_p}\cdots a_{j_1}X
a_{j_1}^\dagger\cdots a_{j_p}^\dagger$, the sum running over all $p$-tuples
of mode indices. A tuple with a repeated index contributes zero, because the
$a_j$ anticommute and $a_j^2=0$. For a tuple of distinct indices with
underlying set $T$, $a_{j_p}\cdots a_{j_1}=\pm A_T$, and the sign cancels
against that of the adjoint. Each set $T$ arises from $p!$ tuples.
\end{proof}

\subsection{The main theorem}

\begin{theorem}\label{thm:main}
Let $\ket\psi\in H_n$ be a unit vector, $\varrho=\ket\psi\bra\psi$ and
$1\le p\le n$. Let $V:\C^{\binom np}\to H_n$ be the linear map defined on the
standard basis $(e_T)_{|T|=p}$ by $Ve_T=A_T\ket\psi$. Then
\[
\Phi^p(\varrho)=p!\,VV^\dagger,\qquad
\bra T\varrho_p\ket S=(V^\dagger V)_{S,T}\quad(|S|=|T|=p),
\]
and the full spectra of the two operators on $H_n$, counted with
multiplicity and zero eigenvalues included, satisfy
$\Spec\Phi^p(\varrho)=p!\,\Spec\varrho_p$. Moreover, if
$c=\sum_Tc_T\ket T$ is an eigenvector of $\varrho_p$ with eigenvalue
$\lambda>0$, then $\sum_T\overline{c_T}\,A_T\ket\psi$ is a nonzero
eigenvector of $\Phi^p(\varrho)$ with eigenvalue $p!\,\lambda$.
\end{theorem}

\begin{proof}
Lemma~\ref{lem:phip} gives
$\Phi^p(\varrho)=p!\sum_TA_T\ket\psi\bra\psi A_T^\dagger=p!\,VV^\dagger$.
Lemma~\ref{lem:units} gives
$\bra T\varrho_p\ket S=\bra\psi A_S^\dagger A_T\ket\psi
=\braket{A_S\psi}{A_T\psi}=(V^\dagger V)_{S,T}$. In the basis $(\ket T)$ of
the $p$-particle sector, the matrix of $\varrho_p$ is therefore the transpose
of the Hermitian matrix $V^\dagger V$, that is, its complex conjugate. A
Hermitian matrix and its complex conjugate have the same spectrum, and
$VV^\dagger$ and $V^\dagger V$ have the same nonzero eigenvalues with the same
multiplicities. The nonzero eigenvalues of $\Phi^p(\varrho)$ are thus those
of $\varrho_p$ multiplied by $p!$. Both operators act on the same space of
dimension $2^n$ and have the same rank, so the multiplicities of the
eigenvalue $0$ agree as well. For the last statement, $\varrho_pc=\lambda c$
reads $\overline{V^\dagger V}\,c=\lambda c$, that is,
$V^\dagger V\bar c=\lambda\bar c$. Then
$VV^\dagger(V\bar c)=\lambda V\bar c$ and
$\norm{V\bar c}^2=\lambda\norm c^2>0$.
\end{proof}

The theorem has a transparent structure. The map
$\mathcal N_p(X)=\sum_{|T|=p}A_TXA_T^\dagger$ is completely positive with
Kraus operators $A_T$, and Lemma~\ref{lem:units} identifies $\varrho_p$ with
the output of the complementary map: in the basis $(\ket T)_{|T|=p}$ of the
$p$-particle sector, $\bra T\varrho_p\ket S=\Tr(A_T\varrho A_S^\dagger)$,
which is the matrix element of the complementary output indexed by the pair
of Kraus operators $(A_T,A_S)$. For a pure input, the output of a
completely positive map and the output of its complement share their nonzero
spectrum. The conjecture is therefore an instance of this general principle,
and the factor $p!$ counts the orderings in which $\Phi^p$ removes the same
set of $p$ particles.

\subsection{Consequences for the conjecture}

\begin{corollary}[$p=1$]\label{cor:p1}
For every pure state on $H_n$, $\Spec\Phi(\varrho)=\Spec\varrho_1$. The
conjecture holds for $p=1$ with no assumption on parity or particle number.
\end{corollary}

\begin{corollary}[$p\ge2$]\label{cor:p2}
Let $p\ge2$ and $\varrho=\ket\psi\bra\psi$. The spectra of $\Phi^p(\varrho)$
and $\varrho_p$ coincide if and only if both operators vanish, that is, if
and only if $\ket\psi$ has no component in the sectors with $p$ or more
particles. In particular the conjecture fails for every $p\ge2$, already for
the two-mode state $\ket{11}$ with $p=2$, for which $\Spec\varrho_2=\{1,0,0,0\}$
and $\Spec\Phi^2(\varrho)=\{2,0,0,0\}$.
\end{corollary}

\begin{proof}
A multiset $\sigma$ of nonnegative reals equals $p!\,\sigma$ with $p!>1$ only
if every element of $\sigma$ is zero. By Theorem~\ref{thm:main}, both
operators vanish exactly when $V=0$, that is, when $A_T\ket\psi=0$ for every
$|T|=p$. Writing $\ket\psi=\sum_J\psi_J\ket J$, the vector $A_T\ket\psi$ is
$\sum_{J\supseteq T}\pm\psi_J\ket{J\setminus T}$, and the vectors
$\ket{J\setminus T}$ are distinct for distinct $J\supseteq T$. Hence $V=0$ if
and only if $\psi_J=0$ whenever $|J|\ge p$. For $\ket{11}$ we have
$\varrho_2=\ket{11}\bra{11}$ by Lemma~\ref{lem:units}, and
$\Phi^2(\varrho)=2A_{\{1,2\}}\ket{11}\bra{11}A_{\{1,2\}}^\dagger
=2\ket{00}\bra{00}$ by Lemma~\ref{lem:phip}.
\end{proof}

The obstruction is already visible on the traces. Since
$\Tr\varrho_p=\sum_{|S|=p}\langle A_S^\dagger A_S\rangle$ and
$\sum_{|S|=p}A_S^\dagger A_S=\binom{\Nop}{p}$, while
$\Tr\Phi^p(\varrho)=\sum_{(j_1,\dots,j_p)}\langle a_{j_1}^\dagger\cdots
a_{j_p}^\dagger a_{j_p}\cdots a_{j_1}\rangle$ counts ordered tuples of
distinct occupied modes,
\[
\Tr\Phi^p(\varrho)=\big\langle\Nop(\Nop-1)\cdots(\Nop-p+1)\big\rangle
=p!\,\Big\langle\binom{\Nop}{p}\Big\rangle=p!\,\Tr\varrho_p
\]
for every state, pure or mixed.

The verdict depends on what one calls the $p$-particle reduced density
matrix. The abstract of \cite{AmosovFilippov2017} speaks of the
conventional $p$-particle reduced density matrix, and the source also writes
it as the kernel indexed by ordered tuples. That reading is equispectral with
$\Phi^p(\varrho)$ on the nonzero part.

\begin{corollary}[tuple-indexed kernel]\label{cor:kernel}
Let $D^{(p)}$ be the $n^p\times n^p$ matrix indexed by ordered $p$-tuples,
$D^{(p)}_{s,t}=\langle a_{s_1}^\dagger\cdots a_{s_p}^\dagger
a_{t_p}\cdots a_{t_1}\rangle$. For every pure state, $D^{(p)}$ and
$\Phi^p(\varrho)$ have the same nonzero eigenvalues with the same
multiplicities.
\end{corollary}

\begin{proof}
Let $W$ map the basis vector $e_{(t_1,\dots,t_p)}$ of $\C^{n^p}$ to
$a_{t_p}\cdots a_{t_1}\ket\psi$. Then $D^{(p)}=W^\dagger W$, and the first
line of the proof of Lemma~\ref{lem:phip} reads $\Phi^p(\varrho)=WW^\dagger$.
\end{proof}

The multiplicities of the eigenvalue $0$ differ, since $n^p\ne2^n$ in
general. The conjecture is therefore exact on nonzero eigenvalues if the
reduced density matrix is the tuple-indexed kernel, and false by the factor
$p!$ for the operator $\varrho_p$ of Definition~\ref{def:rdm}, which is the
operator for which the conjecture is stated in the body of
\cite{AmosovFilippov2017}. In both readings the mathematical content is the
identity \eqref{eq:main}.

\subsection{Fixed particle number}
\begin{corollary}[correction of Proposition~4 of \cite{AmosovFilippov2017}]
\label{cor:fixedN}
If $\varrho$, pure or mixed, is supported on the $N$-particle sector and
$0\le p\le N$, then $\Phi^{N-p}(\varrho)=(N-p)!\,\varrho_p$, where
$\varrho_0=\Tr(\varrho)\ket0\bra0$.
\end{corollary}

\begin{proof}
By linearity it suffices to treat $\varrho=\ket\psi\bra\psi$ with $\ket\psi$
in the $N$-particle sector. Put $m=N-p$. By Lemma~\ref{lem:phip} it suffices
to prove $\sum_{|T|=m}A_T\varrho A_T^\dagger=\varrho_p$; for $m=0$ this reads
$\varrho=\varrho_N$, which is Lemma~\ref{lem:units}, and for $p=0$ both sides
equal $\ket0\bra0$. Both sides are supported on the $p$-particle sector. For
$|J|=|K|=p$, Lemma~\ref{lem:units} and the completeness of the basis
$\{\ket T:|T|=m\}$ of the $m$-particle sector, which contains $A_J\ket\psi$
and $A_K\ket\psi$, give
\[
\bra J\varrho_p\ket K=\braket{A_K\psi}{A_J\psi}
=\sum_{|T|=m}\braket{A_K\psi}{T}\braket{T}{A_J\psi}.
\]
Now $\braket{T}{A_J\psi}=\bra0A_TA_J\ket\psi$ and
$\braket{J}{A_T\psi}=\bra0A_JA_T\ket\psi$, and
$A_JA_T=(-1)^{pm}A_TA_J$ because annihilation operators anticommute. Hence
$\braket{T}{A_J\psi}=(-1)^{pm}\braket{J}{A_T\psi}$, likewise for $K$, and the
two signs cancel:
$\bra J\varrho_p\ket K=\sum_T\braket{J}{A_T\psi}\braket{A_T\psi}{K}
=\bra J\big(\sum_TA_T\varrho A_T^\dagger\big)\ket K$.
\end{proof}

Proposition~4 of \cite{AmosovFilippov2017} asserts
$\Phi^{N-p}(\varrho)=\varrho_p$; it is exact for $N-p\le1$ and off by the
factor $(N-p)!$ beyond. The corollary following it, equality of the spectra
of $\varrho_p$ and $\Phi^p(\varrho)$ for pure fixed-$N$ states, is therefore
also exact only for $p\le1$, as the example $\ket{11}$ shows. Combined with
Theorem~\ref{thm:main}, Corollary~\ref{cor:fixedN} recovers the
Carlson--Keller duality in the normalization of the source: for a pure
$N$-particle state,
$\Spec\varrho_{N-p}=\frac1{p!}\Spec\Phi^p(\varrho)=\Spec\varrho_p$.
Theorem~\ref{thm:main} is thus the extension of this duality to the whole
Fock space, with $\Phi^p(\varrho)/p!$ playing the role of the
$(N-p)$-particle reduced density operator.

\subsection{Mixed states}
The spectral identity does not survive mixing, even for $p=1$ and fixed
particle number. For two modes and
$\varrho=\frac12(\ket{10}\bra{10}+\ket{01}\bra{01})$, the operator
$\varrho_1$ has spectrum $\{\frac12,\frac12,0,0\}$ while
$\Phi(\varrho)=\ket{00}\bra{00}$ has spectrum $\{1,0,0,0\}$. For mixed states
only the trace identity $\Tr\Phi^p(\varrho)=p!\Tr\varrho_p$ and the
complementarity relation survive in general. The source states the
conjecture for pure states only.

\section{Verification}\label{sec:verification}
The
\href{https://github.com/jpbruneton/open_quantum_problems/blob/main/research/entries/fermionic-particle-reductions/checks/verify_particle_reductions.py}{Python checker}
builds the ladder operators with the sign convention of
\cite{AmosovFilippov2017}, implements $\varrho_p$ literally from
Definition~\ref{def:rdm}, projector product included, and $\Phi^p$ by
composition, and writes its
\href{https://github.com/jpbruneton/open_quantum_problems/blob/main/research/entries/fermionic-particle-reductions/data/verification.json}{recorded output}.
It checks Lemma~\ref{lem:units} for $n\le4$; Lemma~\ref{lem:phip} on random
mixed operators for $n\le5$; the identity $\Spec\Phi^p(\varrho)=p!\Spec\varrho_p$
on all $2^n$ eigenvalues for 106 random pure states, real and complex,
covering generic Fock-space states, even- and odd-parity states and
fixed-$N$ states with $N=2,3$ on four to six modes, with largest deviation of
order $10^{-15}$; the eigenvector map of Theorem~\ref{thm:main};
Corollary~\ref{cor:kernel}; Corollary~\ref{cor:fixedN} for pure and mixed
states; the trace identities; the two exact counterexamples above; and the
two-mode example of the source. For $p\ge2$ the largest observed gap between
the spectra of $\Phi^p(\varrho)$ and $\varrho_p$ is about $5.29$, while for
$p=1$ it is below $10^{-15}$. The checker needs only NumPy and runs in a few
seconds.

\section{Scope and limitations}\label{sec:scope}
The analytic proof is complete and elementary. It covers every pure state on
the Fock space of $n$ modes, for all $n$ and $p$, and uses only the canonical
anticommutation relations and the equality of the nonzero spectra of
$VV^\dagger$ and $V^\dagger V$. The numerical computations serve to confirm
the sign and normalization conventions of the source; they are restricted to
$n\le6$ and are not a formal verification.

On bibliographic priority, the factorial correction at fixed particle number
is not new: it appears explicitly, without comment on Proposition~4 of the
source, in the definition of the $M$-body density operator of Gigena, Di
Tullio and Rossignoli \cite[Eq.~(79)]{Gigena2021}, and the fixed-$N$ duality
is classical. The extension to variable particle number, the
identification of $\varrho_p$ and $\Phi^p(\varrho)/p!$ as complementary
outputs, and the observation that Proposition~4 and its corollary in
\cite{AmosovFilippov2017} fail for $N-p\ge2$ and $p\ge2$ were not found in
our search, which covered the 42 works citing \cite{AmosovFilippov2017}
listed by Semantic Scholar in November 2026, with the citation contexts of
the eight closest ones read at source, and keyword web searches. The result
is elementary enough that it may appear elsewhere in another form; the search
did not cover Google Scholar or quantum-chemistry monographs.

The generalized Pauli constraints for states obeying the parity
superselection rule, raised in \cite{AmosovFilippov2017}, remain outside the
scope of this note, as does the spectrum of $\Phi^p(\varrho)$ for mixed
states, where only the trace identity and the complementarity relation are
available.

\begin{thebibliography}{9}
\bibitem{AmosovFilippov2017}
G.~G. Amosov and S.~N. Filippov,
\emph{Spectral properties of reduced fermionic density operators and parity
superselection rule},
Quantum Inf. Process. \textbf{16}, 2 (2017),
\href{https://doi.org/10.1007/s11128-016-1467-9}{doi:10.1007/s11128-016-1467-9},
\href{https://arxiv.org/abs/1512.01828}{arXiv:1512.01828}.

\bibitem{Gigena2021}
N.~Gigena, M.~Di~Tullio and R.~Rossignoli,
\emph{Many-body entanglement in fermion systems},
Phys. Rev. A \textbf{103}, 052424 (2021),
\href{https://doi.org/10.1103/PhysRevA.103.052424}{doi:10.1103/PhysRevA.103.052424},
\href{https://arxiv.org/abs/2012.13785}{arXiv:2012.13785}.
\end{thebibliography}
\end{document}
